% !PRINTED 392
\weberpage{392}

%<<CONT
into such a system $S'$, and we call such a system
\gesperrt{real} when $S$ is identical with $S'$. A real seven-system
therefore contains, for every root $\xi$ occurring in it, the conjugate
$\xi'$, and consequently must contain at least one, but always an
odd number, of real double tangents.

We now derive the principal theorems in order:

\gesperrt{1. Not all the double tangents can be imaginary.}

For let $\xi, \xi'$ be a conjugate pair and $\eta$ a third
imaginary double tangent, such that the three $\xi, \xi', \eta$ are \gesperrt{azygetic}.
(This would certainly be possible if all the double tangents
were imaginary.) If $\eta'$ is conjugate to $\eta$, we have the two
azygetic triples:

\[
\xi,\ \xi',\ \eta;\quad \xi,\ \xi',\ \eta'.
\]

Since, accordingly, $\xi'$ does not occur in the \gesperrt{Steiner}
complex determined by the pair $\xi\eta$, this complex is
imaginary.

We set it beside its conjugate complex:

\weq{1}{
\begin{aligned}
&1)\quad \xi\eta,\ \xi_1\eta_1,\ \xi_2\eta_2,\ \xi_3\eta_3,\ \xi_4\eta_4,\ \xi_5\eta_5\\
&2)\quad \xi'\eta',\ \xi_1'\eta_1',\ \xi_2'\eta_2',\ \xi_3'\eta_3',\ \xi_4'\eta_4',\ \xi_5'\eta_5',
\end{aligned}
}

where by $\xi_i', \eta_i'$ we understand the elements conjugate to $\xi_i, \eta_i$,
so that, if $\xi_i$ is real, we are to set $\xi_i = \xi_i'$, and conversely
from $\xi_i = \xi_i'$ it also follows that $\xi_i$ is real.

Two possibilities are now to be distinguished. If,
\gesperrt{first}, the complexes 1), 2) form a syzygetic pair (\S.\,97),
then they have four syzygetic elements in common. Among these
$\xi, \eta, \xi', \eta'$ cannot occur, and we therefore do not restrict
generality if we assume that $\xi_1, \eta_1, \xi_2, \eta_2$ are
the common elements. But these cannot in their totality
be different from $\xi_1', \eta_1', \xi_2', \eta_2'$, since the
passage to the accented letters, i.\,e. to the conjugate
imaginary quantities, by which 1) goes over into 2), is everywhere
permitted.

Now the assumption $\xi_1 = \eta_1'$ is excluded, because from it
$\xi_1' = \eta_1$ would follow and 1) would not be different from 2).

If $\xi_1 = \xi_1'$, then $\xi_1$ is \gesperrt{real}.

But if $\xi_1 = \xi_2'$, $\xi_2 = \xi_1'$,
then $\eta_1 = \eta_1'$, hence $\eta_1$ is \gesperrt{real},
and from this the assumption $\xi_1 = \eta_2'$, $\xi_1' = \eta_2$, $\xi_2 = \xi_2'$ is not
essentially different.
